% 编译提示：使用 XeLaTeX 编译
% 需要 ctex 宏包支持中文

\documentclass[a4paper,9pt]{article}
\usepackage{amsmath}
\usepackage{amssymb}
\usepackage{geometry}
\usepackage{ctex}
\usepackage{xcolor}
\usepackage{multicol}
\usepackage{titlesec}
\usepackage{fancyhdr}
\usepackage{tcolorbox}
\usepackage{graphicx}
\usepackage[version=4]{mhchem}

\geometry{left=1cm,right=1cm,top=1cm,bottom=1.8cm}
\setlength{\columnsep}{0.8cm}
\setlength{\columnseprule}{0.4pt}
\renewcommand{\columnseprulecolor}{\color{gray!50}}

\definecolor{sectioncolor}{RGB}{0,102,204}
\definecolor{subsectioncolor}{RGB}{0,153,153}
\definecolor{keycolor}{RGB}{204,102,0}
\definecolor{derivcolor}{RGB}{100,149,237}
\definecolor{boxbg}{RGB}{240,248,255}
\definecolor{keybg}{RGB}{255,248,240}

\titleformat{\section}{\normalfont\large\bfseries\color{sectioncolor}}{\thesection}{0.5em}{}
\titlespacing*{\section}{0pt}{1ex plus 0.5ex minus 0.2ex}{0.8ex plus 0.2ex}
\titleformat{\subsection}{\normalfont\normalsize\bfseries\color{subsectioncolor}}{\thesubsection}{0.5em}{}
\titlespacing*{\subsection}{0pt}{0.8ex plus 0.3ex minus 0.1ex}{0.5ex plus 0.1ex}
\setcounter{secnumdepth}{4}\setcounter{tocdepth}{4}

\newenvironment{keypoint}
{\par\vspace{0.3em}\noindent\begin{tcolorbox}[colback=keybg,colframe=keycolor,boxrule=0.5pt,arc=2pt,left=2pt,right=2pt,top=2pt,bottom=2pt,boxsep=0pt]\noindent\textcolor{keycolor}{\textbf{关键点：}}}
{\end{tcolorbox}\par\vspace{0.3em}}

\newenvironment{examplebox}
{\par\vspace{0.3em}\noindent\begin{tcolorbox}[colback=boxbg,colframe=derivcolor,boxrule=0.5pt,arc=2pt,left=2pt,right=2pt,top=2pt,bottom=2pt,boxsep=0pt]\scriptsize\noindent\textcolor{derivcolor}{\textbf{示例：}}\par\setlength{\parskip}{0.1ex}\setlength{\itemsep}{0pt}\setlength{\parsep}{0pt}}
{\end{tcolorbox}\par\vspace{0.3em}}

\pagestyle{fancy}
\fancyhf{}
\fancyfoot[C]{\scriptsize11-\thepage}
\renewcommand{\headrulewidth}{0pt}
\renewcommand{\footrulewidth}{0.4pt}
\renewcommand{\footrule}{\hbox to\headwidth{\color{gray!50}\leaders\hrule height \footrulewidth\hfill}}

\title{\vspace{-2em}\Large\bfseries\color{sectioncolor} Chapter 11 Collections - Eyring方程、活化参量与Boltzmann分布（Lect.12）\vspace{-1em}}
\author{}
\date{\today}

\begin{document}
\maketitle
\thispagestyle{fancy}

\begin{multicols}{2}
\scriptsize{\tableofcontents}

\section{过渡态理论的热力学表述}

由 lect11 的结果
\[
k_{2\mathrm{nd}}=\left(\frac{1}{c^\circ}\right)\frac{kT}{h}K^{\ddagger}
\]
出发，再用
\[
\Delta_rG^{\circ,\ddagger}=-RT\ln K^{\ddagger},
\qquad
\Delta_rG^{\circ,\ddagger}=\Delta_rH^{\circ,\ddagger}-T\Delta_rS^{\circ,\ddagger},
\]
即可把速率常数写成热力学活化参数的形式。对溶液中反应：
\[
k_{2\mathrm{nd}}=\frac{1}{[\ ]^\circ}\frac{kT}{h}
\exp\left(\frac{\Delta_rS^{\circ,\ddagger}}{R}\right)
\exp\left(-\frac{\Delta_rH^{\circ,\ddagger}}{RT}\right).
\]

对气相，则需要改用以标准压强为基准的表达式，标准态不同，单位也不同。

\begin{keypoint}
Eyring 方程的价值不在“更容易算”，而在于它把速率常数拆成了\textbf{活化焓项}和\textbf{活化熵项}，因此非常适合解释机理与环境效应。
同时要记住：
\[
\Delta_rG^{\circ,\ddagger}=-RT\ln K^{\ddagger}
\]
只有在 $K^{\ddagger}$ 与 $\Delta_rG^{\circ,\ddagger}$ 使用\textbf{同一个标准态}时才能直接成立。
\end{keypoint}

\begin{keypoint}
	本讲建议固定下面这套单位习惯：$\Delta_rG^{\circ,\ddagger}$、$\Delta_rH^{\circ,\ddagger}$、$E_a$ 用 $\mathrm{kJ\,mol^{-1}}$ 或 $\mathrm{J\,mol^{-1}}$，并始终与 $R=8.314\ \mathrm{J\,mol^{-1}\,K^{-1}}$ 保持一致；$\Delta_rS^{\circ,\ddagger}$ 用 $\mathrm{J\,mol^{-1}\,K^{-1}}$。若把焓写成 kJ/mol 而仍直接代入以 J 为单位的 $R$，数值会差 1000 倍。
\end{keypoint}

\section{与 Arrhenius 方程的关系}

Arrhenius 方程为
\[
k=A\exp\left(-\frac{E_a}{RT}\right).
\]
但不能简单把 $E_a$ 认成 $\Delta_rH^{\circ,\ddagger}$，因为活化能的正式定义是
\[
\left(\frac{\partial \ln k}{\partial (1/T)}\right)_V=-\frac{E_a}{R},
\]
或等价地
\[
\left(\frac{\partial \ln k}{\partial T}\right)_V=\frac{E_a}{RT^2}.
\]

对溶液反应，结合 van't Hoff isochore 得
\[
E_a=\Delta_rH^{\circ,\ddagger}+RT.
\]
对\textbf{双分子气相反应}，因 $K_p$ 和 $K_c$ 的关系中多出体积/压强项，结果是
\[
E_a=\Delta_rH^{\circ,\ddagger}+2RT.
\]

\begin{keypoint}
\textbf{考试最容易错的点之一}：
\[
E_a \neq \Delta_rH^{\circ,\ddagger}
\]
且修正项依赖于所讨论的是溶液反应还是双分子气相反应。
\end{keypoint}

\section{Eyring plot 与活化参量实验求取}

把 Eyring 方程改写为
\[
\ln\left(\frac{k_{2\mathrm{nd}}}{T}\right)
=\ln\left(\frac{k}{h[\ ]^\circ}\right)
+\frac{\Delta_rS^{\circ,\ddagger}}{R}
-\frac{\Delta_rH^{\circ,\ddagger}}{R}\frac{1}{T}.
\]
故作图 $\ln(k/T)$ 对 $1/T$ 会得到直线：
\begin{itemize}
  \item 斜率 $=-\Delta_rH^{\circ,\ddagger}/R$；
  \item 截距 $=\ln(k/h[\ ]^\circ)+\Delta_rS^{\circ,\ddagger}/R$。
\end{itemize}
这里隐含的实验拟合假设是：在所测温区内，$\Delta_rH^{\circ,\ddagger}$ 和 $\Delta_rS^{\circ,\ddagger}$ 可近似视为与温度无关。

也可以先做 Arrhenius plot 求 $E_a$ 与 $A$，再回算活化焓与活化熵，但由于实验常只覆盖有限温区，截距外推会带来较大误差，因此活化熵往往不如活化焓稳健。

\section{活化体积与机理判断}

由主方程
\[
dG=Vdp-SdT
\]
在恒温下得
\[
\left(\frac{\partial G}{\partial p}\right)_T=V.
\]
因此
\[
\Delta_rV^{\circ,\ddagger}=\left(\frac{\partial \Delta_rG^{\circ,\ddagger}}{\partial p}\right)_T,
\]
并且
\[
\left(\frac{\partial \ln k_{2\mathrm{nd}}}{\partial p}\right)_T
=-\frac{\Delta_rV^{\circ,\ddagger}}{RT}.
\]
这一套讨论主要是针对\textbf{溶液反应}在高压下的速率常数变化；普通气相反应的压强效应往往先体现为浓度变化，而不是这里讨论的活化体积效应。

若 $\Delta_rV^{\circ,\ddagger}$ 与压强近似无关，则
\[
\ln\frac{k_{2\mathrm{nd}}(p)}{k_{2\mathrm{nd}}(p^\circ)}
=-\frac{\Delta_rV^{\circ,\ddagger}}{RT}(p-p^\circ).
\]

\begin{keypoint}
	活化体积公式里，$\Delta_rV^{\circ,\ddagger}(p-p^\circ)$ 必须先组成“每摩尔能量”再去除以 $RT$。最稳妥的做法是统一到 SI：$\Delta_rV^{\circ,\ddagger}$ 用 $\mathrm{m^3\,mol^{-1}}$，$p$ 用 Pa；因为 $\mathrm{Pa\cdot m^3=J}$，这样就能直接与 $RT$ 的 $\mathrm{J\,mol^{-1}}$ 对上。
\end{keypoint}

\begin{examplebox}
机理判断的经验规则：
\begin{itemize}
  \item 缔合型机理：$\Delta_rS^{\circ,\ddagger}$、$\Delta_rV^{\circ,\ddagger}$ 常偏负；
  \item 解离型机理：二者常偏正；
  \item 互换型机理：$\Delta_rS^{\circ,\ddagger}$ 往往接近 0，$\Delta_rV^{\circ,\ddagger}$ 可正可负。
\end{itemize}
lect12 特别强调：\textbf{活化体积通常比活化熵更适合判机理。}
\end{examplebox}

\begin{examplebox}
电致收缩（electrostriction）是 lect12 解释溶液机理的核心词：离子被极性溶剂强烈配位时，溶剂熵与体积都会下降。若过渡态中\textbf{电荷分散}，电致收缩减弱，则常有
\[
\Delta_rS^{\circ,\ddagger}>0,\qquad \Delta_rV^{\circ,\ddagger}>0.
\]
若过渡态中\textbf{电荷分离}增强，则更可能出现
\[
\Delta_rV^{\circ,\ddagger}<0.
\]
这正是为什么这两类活化参量可以拿来区分缔合、解离和电荷重排机理。
\end{examplebox}

\section{Boltzmann 分布的重新推导}

对非相互作用粒子，单粒子配分函数为
\[
q=\sum_i e^{-\varepsilon_i/kT}.
\]
内能可由配分函数写成
\[
U=NkT^2\left(\frac{\partial \ln q}{\partial T}\right)_V.
\]
而
\[
\left(\frac{\partial q}{\partial T}\right)_V=\frac{1}{kT^2}\sum_i \varepsilon_i e^{-\varepsilon_i/kT},
\]
故
\[
U=\frac{N}{q}\sum_i \varepsilon_i e^{-\varepsilon_i/kT}.
\]
另一方面，若按各态布居数求和，
\[
U=\sum_i n_i\varepsilon_i.
\]
逐项比较得到
\[
n_i=\frac{N}{q}e^{-\varepsilon_i/kT}.
\]
若按能级求和且有简并度 $g_j$，则
\[
n_j=\frac{Ng_j}{q}e^{-\varepsilon_j/kT}.
\]

\begin{keypoint}
lect12 给出的一个很实用解释是：若基态取零点，则
\[
n_0=\frac{N}{q},
\qquad q=\frac{N}{n_0}.
\]
所以 $q$ 可看成“离开基态的程度”的量度；越多粒子爬离基态，$q$ 越大。
\end{keypoint}

\section{Boltzmann 分布的应用}

\subsection{振动布居}
若简谐振子能级取为
\[
E_v=vh\nu_0,
\qquad
q'_{\mathrm{vib}}=\frac{1}{1-e^{-\theta_{\mathrm{vib}}/T}},
\]
则振动态布居为
\[
n_v=N\left[1-e^{-\theta_{\mathrm{vib}}/T}\right]e^{-v\theta_{\mathrm{vib}}/T}.
\]

\subsection{转动布居与 $J_{\max}$}
刚性转子中
\[
E_J=BJ(J+1),
\qquad
n_J=\frac{N(2J+1)}{q_{\mathrm{rot}}}e^{-BJ(J+1)/kT}.
\]
当 $B\ll kT$ 时，
\[
q_{\mathrm{rot}}=\frac{T}{\sigma\theta_{\mathrm{rot}}},
\qquad
\theta_{\mathrm{rot}}=\frac{B}{k},
\]
故
\[
\frac{n_J}{N}=\frac{\sigma\theta_{\mathrm{rot}}}{T}(2J+1)
e^{-\theta_{\mathrm{rot}}J(J+1)/T}.
\]
对 $J$ 求导并令其为 0，得最可几转动量子数：
\[
J_{\max}=\sqrt{\frac{T}{2\theta_{\mathrm{rot}}}}-\frac{1}{2}.
\]

\begin{examplebox}
CO 在 298 K 的课件例子中，$J_{\max}\approx 6.8$，所以真正整数最大值在
\[
J=7.
\]
这也是解释双原子分子红外光谱 P/R 支强度分布的关键。
\end{examplebox}

\section{非平衡态与负温度}

平衡时，作图 $\ln n_i$ 对 $\varepsilon_i$ 应为直线，斜率
\[
\text{slope}=-\frac{1}{kT}.
\]
若体系被强辐射或快速反应扰乱后偏离平衡，则这一斜率定义的“有效温度”也会改变。对两能级体系，
\[
\frac{n_1}{n_0}=\exp\left(-\frac{\Delta\varepsilon}{kT}\right).
\]
热平衡时总有 $n_0>n_1$。若人为造成
\[
n_1>n_0,
\]
则形式上对应
\[
T<0,
\]
这就是\textbf{population inversion}，也就是负温度态。

\section{Bose--Einstein 与 Fermi--Dirac 统计简介}

前面整套笔记主要默认的是不发生量子简并的 Maxwell--Boltzmann 统计。若粒子不可分辨且量子态占据受量子规则限制，则平均占据数需要改写成
\[
\bar n_i^{\mathrm{BE}}=\frac{g_i}{\exp[(\varepsilon_i-\mu)/kT]-1},
\]
\[
\bar n_i^{\mathrm{FD}}=\frac{g_i}{\exp[(\varepsilon_i-\mu)/kT]+1}.
\]
前者适用于 boson，可多个粒子占据同一态；后者适用于 fermion，每个单粒子态最多占 1 个粒子。

\begin{keypoint}
当
\[
\exp[(\varepsilon_i-\mu)/kT]\gg 1,
\]
两式都退化为 Maxwell--Boltzmann 形式
\[
\bar n_i\approx g_i e^{-(\varepsilon_i-\mu)/kT}.
\]
这就是为什么常温常压下的大多数分子气体仍然可以安全使用 Boltzmann 分布。
\end{keypoint}

\section{态密度}

态密度定义为
\[
D(\varepsilon)d\varepsilon=W(\varepsilon+d\varepsilon)-W(\varepsilon),
\qquad
D(\varepsilon)=\frac{dW(\varepsilon)}{d\varepsilon},
\]
其中 $W(\varepsilon)$ 是能量不超过 $\varepsilon$ 的状态总数。

对三维立方箱中粒子，能级为
\[
E_{n_x,n_y,n_z}=(n_x^2+n_y^2+n_z^2)\frac{h^2}{8ma^2}.
\]
把每个态看成量子数空间中的一个格点，则满足 $E\le \varepsilon$ 的态都落在半径 $n_{\max}$ 的球内，其中
\[
n_{\max}=\left(\frac{8ma^2\varepsilon}{h^2}\right)^{1/2}.
\]
由于 $n_x,n_y,n_z$ 都是正整数，只需取球的正八分体体积，因此
\[
W(\varepsilon)=\frac{\pi}{6}\left(\frac{8m\varepsilon}{h^2}\right)^{3/2}V.
\]
进一步微分得
\[
D(\varepsilon)=\frac{\pi}{4}\left(\frac{8m}{h^2}\right)^{3/2}\varepsilon^{1/2}V.
\]

\begin{keypoint}
态密度的本质作用是把“极密的平动能级”由离散求和转换成近似连续描述。对理想气体平动来说，正是因为 $D(\varepsilon)$ 很大，常温下才会显得近乎经典。
\end{keypoint}

\section{最后速记}

\begin{keypoint}
lect12 高频公式：
\[
k_{2\mathrm{nd}}=\left(\frac{1}{c^\circ}\right)\frac{kT}{h}K^{\ddagger},
\qquad
\Delta_rG^{\circ,\ddagger}=\Delta_rH^{\circ,\ddagger}-T\Delta_rS^{\circ,\ddagger},
\]
\[
k_{2\mathrm{nd}}=\frac{1}{[\ ]^\circ}\frac{kT}{h}
\exp\left(\frac{\Delta_rS^{\circ,\ddagger}}{R}\right)
\exp\left(-\frac{\Delta_rH^{\circ,\ddagger}}{RT}\right),
\]
\[
\begin{aligned}
E_a&=\Delta_rH^{\circ,\ddagger}+RT && (\text{solution}),\\
E_a&=\Delta_rH^{\circ,\ddagger}+2RT && (\text{bimolecular gas phase}).
\end{aligned}
\]
\[
\left(\frac{\partial \ln k}{\partial p}\right)_T=-\frac{\Delta_rV^{\circ,\ddagger}}{RT},
\qquad
n_i=\frac{N}{q}e^{-\varepsilon_i/kT},
\qquad
n_j=\frac{Ng_j}{q}e^{-\varepsilon_j/kT},
\]
\[
\begin{aligned}
J_{\max}&=\sqrt{\frac{T}{2\theta_{\mathrm{rot}}}}-\frac{1}{2},\\
\bar n_i^{\mathrm{BE}}&=\frac{g_i}{e^{(\varepsilon_i-\mu)/kT}-1},\\
\bar n_i^{\mathrm{FD}}&=\frac{g_i}{e^{(\varepsilon_i-\mu)/kT}+1}.
\end{aligned}
\]
\[
D(\varepsilon)=\frac{dW}{d\varepsilon},
\qquad
W(\varepsilon)=\frac{\pi}{6}\left(\frac{8m\varepsilon}{h^2}\right)^{3/2}V.
\]
\end{keypoint}

\end{multicols}
\end{document}
